\chapter*{Abstract}
\label{sec:abstract}
\begin{adjustwidth}{4cm}{4cm}

\par\medskip\noindent\textbf{Problem}\par\nobreak%
Logic synthesis restructures a circuit's And-Inverter Graph (AIG) representation for area, power, and delay, but assembling an effective sequence of transformations remains a heuristic trial-and-error process. Predicting a circuit's optimizability beforehand, using a Graph Neural Network (GNN) trained directly on its structure, could replace that search with a data-driven choice. Industrial AIGs reach millions of nodes, however, exceeding the memory of a single device.

\par\medskip\noindent\textbf{Gap}\par\nobreak%
Graph reduction (partitioning, sparsification, and summarization) is the standard remedy for graphs that exceed device memory, well studied on citation and social networks but never systematically evaluated on AIGs. Cutting an edge there can sever a causal cone and destroy the structural signal a predictor depends on.

\par\medskip\noindent\textbf{Approach}\par\nobreak%
This thesis benchmarks generic and AIG-adapted reductions from all three families for optimizability regression, weighing predictive retention against efficiency relative to an unreduced baseline. One summarization variant merges only nodes a bounded-depth encoder cannot distinguish, a coarsening with zero accuracy cost by construction. A further test asks whether a model trained solely on reduced graphs still predicts accurately on full, unreduced circuits.

\par\medskip\noindent\textbf{Findings}\par\nobreak%
\todo{State the RQ2/RQ3 headline trade-off (memory reduction at matched predictive retention) and the RQ5 transfer result once measured.}

\end{adjustwidth}

\vspace{1em}
\noindent\textbf{Keywords:}
electronic design automation; logic synthesis; graph neural network; And-Inverter Graph; graph reduction; graph coarsening
